\documentclass[conference]{IEEEtran}

\usepackage{iftex}
\ifPDFTeX
\usepackage[utf8]{inputenc}
\usepackage[T1]{fontenc}
\fi
\usepackage{amsthm}
\usepackage{amsmath,bm}
\usepackage{newtxtext,newtxmath}
\usepackage{textcomp}
\interdisplaylinepenalty=2500
\usepackage{xcolor}
\usepackage{graphicx}
\usepackage{array,booktabs,makecell,adjustbox}
\usepackage{mathtools}
\usepackage{enumitem}
\usepackage{siunitx}
\usepackage[final]{microtype}
\usepackage{cite}
\usepackage[breaklinks=true,hidelinks]{hyperref}
\usepackage{bookmark}
\usepackage[nameinlink,capitalise,noabbrev]{cleveref}
\usepackage{pgf}

% Keep this --- needed by PGF outputs that reference \mathdefault
\providecommand{\mathdefault}[1]{#1}

% Generated figures live beside this file only after plots/kalman_carma has
% run. Wrap every include so the manuscript still compiles from a clean
% checkout.
\newcommand{\IncludePGF}[1]{%
  \IfFileExists{#1}{\resizebox{\linewidth}{!}{\input{#1}}}%
                   {\fbox{\parbox{0.92\linewidth}{\centering\small
                     Figure \texttt{\detokenize{#1}} not generated.\\
                     Run \texttt{plots/kalman\_carma/draw\_plots.sh}.}}}%
}

\DeclareRobustCommand{\vct}[1]{\bm{#1}}
\DeclareRobustCommand{\mat}[1]{\bm{#1}}
\newcommand{\R}{\mathbb{R}}
\sisetup{detect-weight=true,detect-family=true,per-mode=symbol}
\emergencystretch=2em
\allowdisplaybreaks

\newtheoremstyle{bodysmall}{3pt}{3pt}{\normalfont\small}{0pt}{\bfseries\small}{.}{0.5em}{}
\theoremstyle{bodysmall}
\newtheorem{proposition}{Proposition}
\newtheorem{remark}{Remark}

\title{Three-Axis CARMA Motion Priors for Marine Inertial Wave Estimation:
An Experimental Comparison with the OU-III Filter (draft)}
\author{%
  \IEEEauthorblockN{Mikhail S. Grushinskiy}
  \IEEEauthorblockA{Independent Researcher, USA, Fair Lawn\\
  Bareboat Necessities GitHub Project\\
  Email: mgrouch@users.github.com}%
}

\begin{document}
\maketitle

\begin{abstract}
This experimental paper replaces the three integrated Ornstein--Uhlenbeck
(OU) translational chains of the shipping OU-III marine inertial estimator by
stationary continuous-time ARMA (CARMA) displacement priors on all three
axes, keeping its attitude, bias, magnetometer, startup and adaptation
unchanged. The priors are CARMA($p$,$p-3$) models with a fitted JONSWAP shape,
moved by one measured frequency scale and discretized exactly. A numerator
zero at DC makes the integral of the displacement stationary, and we show
that OU-III's recurring $S=0$ pseudo-measurement is itself equivalent, for
the smoothing posterior, to a broad DC-zero CARMA(4,1) prior. In a paired
complete-INS simulation benchmark (31 sea conditions, 20--100 seeds each)
CARMA(4,1) on all three axes, at OU-III's 21 states and a third of its
multiply-adds per sample, reduced the horizontal RMS error by 43--73\,\% on
single-peaked, bimodal, broadband, underway and translating seas and by
17\,\% on a mixed swell and wind sea, and the heave error by 27--46\,\%
except on the mixed sea, where it is 13\,\% worse. CARMA(5,2) collapses on
the mixed sea because the frequency mapping places its spectral notch on the
swell. All CARMA priors show more phantom displacement on a still IMU than
OU-III. The results are simulation only; nothing here changes the shipping
estimator or bears on its stability study.
\end{abstract}
\input{../simulation-dataset.tex-part}

\begin{IEEEkeywords}
CARMA process, Kalman filter, inertial navigation, marine heave, horizontal
wave motion, Ornstein--Uhlenbeck process, embedded estimation
\end{IEEEkeywords}

\section{Introduction and Status}
\label{sec:intro}
A strapdown IMU on a small vessel senses the wave-induced acceleration of the
hull together with gravity, sensor bias and attitude error. Displacement is
recovered by double integration, which only a motion prior keeps bounded. The
shipping OU-III filter (\texttt{src/kalman\_ou\_iii}) gives each axis an
integrated OU chain $[v,p,S,a_w]$: the wave acceleration $a_w$ is an OU
process with time constant $\tau$ and stationary standard deviation
$\sigma_{aw}$, $S'=p$, and a pseudo-measurement $S=0$ every $T_S$ with the
SpectralMSE standard deviation $r_S$ bounds the integral of the
displacement. The prior is nearly shape-free above its corner $1/\tau$.

A CARMA prior gives the displacement a rational spectrum with resonant
poles~\cite{brockwell01,brockwell04}. For heave, which a hydrostatic restoring
force makes resonant, that is a natural model; this paper asks whether it is
also a better prior for surge and sway, which have no restoring force, and
what the three-axis filter costs. The contribution is an experimental filter
(\texttt{src/kalman\_carma}), the analysis of the auxiliary state $S$ under a
CARMA prior, and a paired benchmark against the unchanged OU-III.

\emph{Status.} The CARMA filter is a research prototype. It is not included
by the deployed configuration, it has no hardware evidence, and no result of
this paper is a stability, admissibility or hardware statement. The OU-III
stability study concerns the shipping OU-III filter only; a CARMA prior is a
statistical model, not a law imposed on the vessel's physical history. The
full report with every table is \texttt{docs/carma-three-axis.md}.

\section{Three-Axis Wave Motion}
\label{sec:physics}
Axes are earth-fixed NED; the wave coordinate $\vct{p}$ is the oscillatory
displacement about the local equilibrium.

\emph{Heave.} Hydrostatic restoring gives a resonant response amplitude
operator (RAO) times hull averaging $e^{-kd}$; the heave spectrum is the
incident spectrum reshaped by it, single-peaked in single seas and
two-peaked in swell plus wind sea.

\emph{Surge and sway.} A small free-floating hull in long waves follows the
water-particle orbit: in deep water the horizontal displacement along the
propagation direction has the elevation amplitude and lags it by a quarter
period, while hull averaging attenuates short waves. With $b$ the
propagation direction relative to the bow and $k$ the wavenumber, the
benchmark uses
\begin{equation}
H_{\mathrm{surge}}=e^{-k(d+\ell_L|\cos b|)}\cos b,\quad
H_{\mathrm{sway}}=e^{-k(d+\ell_B|\sin b|)}\sin b,
\label{eq:horiz-rao}
\end{equation}
($d=\SI{0.5}{m}$, $\ell_L=\SI{0.75}{m}$, $\ell_B=\SI{0.28}{m}$) times the
quarter-period lag: a smooth low-pass reshaping of the incident spectrum
without resonance.

\emph{Zero at DC.} The elevation spectrum and every RAO vanish at zero
frequency, so the wave displacement \emph{and} its time integral are
bounded. A physically consistent prior therefore has a numerator zero at DC.

\emph{Low-frequency translation.} Anchored sailing on the rode (periods of
\SIrange{30}{300}{s}), wind drift and manoeuvres move the local equilibrium,
not the wave coordinate. Their accelerations (1--30\,mg) are of the size of
accelerometer bias and of $g\,\delta\theta$; an IMU-only estimator cannot
reconstruct them, and the relevant question is how much of them each prior
turns into phantom wave displacement.

\emph{Cross-axis coupling.} For linear waves the cross-spectral matrix of
$(x_N,y_E,z)$ is $S_{xyz}(\omega)=\sum_jS_\eta(\omega)w_j\vct{t}_j
\vct{t}_j^{\mathsf H}$, Hermitian positive semidefinite. Horizontal axes are
mutually in phase and in quadrature with the incident elevation; the vessel's
heave response adds its own phase.

\section{CARMA Priors}
\label{sec:carma}
A CARMA($p$,$q$) displacement with transfer function $gB(s)/P(s)$ ($B$, $P$
monic) driven by unit white noise is realized in the coordinates
\begin{equation}
\vct{x}=[p,\,v,\,a,\,l_1,\dots,l_{N-3}]^{\mathsf T},\qquad
d\vct{x}=\mat{A}\vct{x}\,dt+\vct{G}\,dW,
\label{eq:carma-ss}
\end{equation}
with $\vct{G}=g\,\vct{e}_a$ and $p=\vct{e}_0^{\mathsf T}\vct{x}$. Relative
degree $p-q=3$ gives $\mat{C}\vct{G}=\mat{C}\mat{A}\vct{G}=0$ and
$\mat{C}\mat{A}^2\vct{G}=g$: displacement and velocity are differentiable,
acceleration is continuous with finite variance
($S_a=\omega^4S_p\sim\omega^{-2}$), and its derivative carries the white
noise. The latent coordinates are those of the cascade realization:
CARMA(4,1) has one latent state, the input $r$ of its second oscillator, and
CARMA(5,2) two, $(r,r')$. In these coordinates OU-III's chain is the special
case $a'=-a/\tau+\text{noise}$ without displacement or velocity feedback;
the CARMA feedback of $p$ and $v$ into the jerk equation is what makes the
displacement stationary.

\emph{Shapes.} \texttt{tools/carma\_fit.py} minimizes the RMS log-spectral
error against JONSWAP ($f_p=\SI{0.29}{Hz}$, $\gamma=3.3$) over
\SIrange{0.17}{1.30}{Hz}, shape only, with every pole damping in $[0.05,1]$
(\cref{tab:shapes}). The CARMA(4,1) optimum places its numerator zero at DC.

\begin{table}[t]
\centering
\caption{Library shapes fitted to JONSWAP (Hz, damping). RMS and maximum
log-spectral error over \SIrange{0.17}{1.30}{Hz}.}
\label{tab:shapes}
\footnotesize
\setlength{\tabcolsep}{3pt}
\begin{adjustbox}{max width=\columnwidth}
\begin{tabular}{@{}llccc@{}}
\toprule
Model & Poles and zeros & RMS dB & max dB & DC re peak\\
\midrule
C52J & \makecell[l]{0.278/0.097, 0.443/0.523,\\ real 1.864; zeros 0.151/0.20} & 0.84 & 9.7 & $-23.2$\,dB\\
C41J & 0.307/0.104, 0.815/0.478; zero 0 & 1.86 & 17.5 & zero\\
C30J & 0.386/0.267, real 5.0 (bound) & 2.93 & 24.0 & $-5.8$\,dB\\
\bottomrule
\end{tabular}
\end{adjustbox}
\end{table}

\emph{Scaling.} One shape at $\kappa=1$ is moved by a frequency scale
$\kappa$ (every pole and zero times $\kappa$). With unit gain
$\mathrm{Var}_\kappa(a)=\mathrm{Var}_1(a)/\kappa$ and
$\mathrm{Var}_\kappa(p)=\mathrm{Var}_1(p)/\kappa^5$, so a requested
stationary acceleration standard deviation $\sigma_a$ is met in closed form,
\begin{equation}
g^2=\sigma_a^2\,\kappa/\mathrm{Var}_1(a),
\label{eq:scaling}
\end{equation}
without a run-time Lyapunov solve.

\emph{Exact discretization.} The exact transition~\cite{sarkka23}
$\mat{F}=e^{\mat{A}h}$ and
$\mat{Q}=\int_0^he^{\mat{A}t}\mat{W}e^{\mat{A}^{\mathsf T}t}dt$ are evaluated
by a Taylor series at $h_s=h/2^s$ with $\|\mat{A}h_s\|_\infty\le1/4$ in
diagonally scaled coordinates $\vct{x}=\mat{D}\vct{z}$,
$\mat{D}=\mathrm{diag}(h_s^{d_i})$ ($d_i$ the coupling distance of state $i$
from the noise input), followed by $s$ exact doublings
\begin{equation}
\mat{Q}(2h)=\mat{Q}(h)+\mat{F}\mat{Q}\mat{F}^{\mathsf T},\qquad
\mat{F}(2h)=\mat{F}^2.
\label{eq:doubling}
\end{equation}
The scaling keeps the small entries ($Q_{pp}\sim h^5$) at full relative
precision, and no stationary difference $\mat{P}-\mat{F}\mat{P}\mat{F}^{\mathsf T}$,
which cancels catastrophically in single precision at \SI{5}{ms}, is formed.
The unit tests check the realization against $g^2|B|^2/|P|^2$ ($<10^{-9}$
relative), $\mat{F}$ and $\mat{Q}$ against a double-precision Van Loan
exponential, the semigroup law for non-uniform steps, the stationary
covariance against a Kronecker Lyapunov solve, and single-precision
$\mat{F},\mat{Q}$ within $1.5\times10^{-6}$ of double over
$f_z=\SIrange{0.05}{1}{Hz}$. Long simulated realizations reproduce the
designed variances (within 1\,\%) and spectra (within \SI{0.13}{dB}), with
Brownian acceleration increments.

\section{The Auxiliary State $S$}
\label{sec:s}
OU-III needs $S=\int p\,dt$ and $S=0$ because its displacement is a double
integral of an OU acceleration, so $(v,p)$ alone is not detectable from
acceleration. A CARMA prior is Hurwitz and observable from acceleration; the
question is whether the integral of its displacement is controlled. Append
$S'=p$; $S$ has transfer function $gB(s)/(s\,P(s))$.

\begin{proposition}
If $B(0)\neq0$, the zero eigenvalue of the augmented system is controllable
from the noise and unobservable from acceleration; $S$ is a random walk with
$\mathrm{Var}\,S(t)\sim S_p(0)\,t$. If $B(0)=0$, the zero eigenvalue is
uncontrollable, $S(t)=S_\infty(t)-S_\infty(0)$ with the stationary
$S_\infty=g\,(B(s)/s)/P(s)$ a linear functional of the CARMA state, and
$\mathrm{Var}\,S(t)\to2\,\mathrm{Var}\,S_\infty$. The augmented pair is then
not stabilizable on the unit circle, the discrete algebraic Riccati equation
has no stabilizing solution, and a recurring $S=0$ pseudo-measurement
re-measures $S_\infty$, information the prior already contains.
\end{proposition}

A zero at DC therefore supplies the low-frequency constraint without an
auxiliary state. OU-III reaches the same structure by its pseudo-measurement.

\begin{proposition}
Condition the integrated OU displacement prior
$S_{OU}(\omega)=2\sigma^2\tau/(\omega^4(1+\omega^2\tau^2))$ on $S=0$
observed at information rate $1/N$, $N=r_S^2T_S$. The conditional
displacement spectrum is
\begin{align}
S_{\mathrm{eff}}&=\frac{S_{OU}N\omega^2}{S_{OU}+N\omega^2}
=\frac{2\sigma^2\tau N\omega^2}{2\sigma^2\tau+N\omega^6+N\tau^2\omega^8}
\nonumber\\
&=\frac{2\sigma^2}{\tau}\,\frac{|i\omega|^2}{|P(i\omega)|^2},
\label{eq:ou3-dc-zero}
\end{align}
with $P$ the monic Hurwitz factor of $N\tau^2s^8-Ns^6+2\sigma^2\tau$: a
CARMA(4,1) with $B(s)=s$.
\end{proposition}

The equivalence holds for the two-sided (smoothing) posterior; the causal
filters differ by up to half of the gain in \SIrange{0.1}{1}{Hz}, because
the causal OU-III filter conditions only on past pseudo-measurements. At the
JONSWAP \SI{0.29}{Hz} operating point the factor has two broad pole pairs
near the peak (damping 0.39 and 0.93): OU-III is a broadband DC-zero prior,
whereas the fitted CARMA(4,1) is a narrow one (damping 0.10). OU-III's $S=0$
must therefore not be carried over to a DC-zero CARMA; a CARMA without a DC
zero (the fitted CARMA(5,2)) can carry $S$ with a pole-placement law
$r_S=\sqrt{S_p(0)}\,\tau/(\kappa_R\sqrt{T_S})$, $\kappa_R=0.3627$, as a
regularizer of a mis-specified DC level.

\section{Cross-Axis Coupling}
\label{sec:coupling}
Correlated driving noise through identical horizontal dynamics,
$\mat{Q}_H=\mat{\Sigma}_h\otimes\mat{Q}_1$, is exact and costs no state, but
its cross-spectrum $\Sigma_{xy}|H(i\omega)|^2$ is real: it represents the
in-phase horizontal co-spectrum only. Quadrature coupling to heave needs
different dynamics, e.g.\ a shared latent built on the narrow-band identity
$\xi\approx\omega\int\eta$. With the exact cross-spectrum, the noncausal
Wiener gain of using all three accelerometer axes is 5--7\,\% on
single-peaked seas under white noise (the $x$--$y$ part at most 1.6\,\%).
The causal shared-latent model, even with oracle coupling coefficients, is
worse than independent axes on those seas, because the errors are dominated
by spectral mismatch, which the narrow-band quadrature adds to. The filter
therefore uses independent axes; one horizontal model serves $x$ and $y$.

\section{Filter and Implementation}
\label{sec:filter}
\texttt{Kalman3D\_Wave\_CARMA3<T,NHC,SH,NZC,SZ>} carries the nine attitude,
gyro-bias and accelerometer-bias error states of OU-III and a CARMA block
$[p,v,a,\text{latent}]$ per axis (orders $3$ to $5$), with optional
$S_x,S_y$ ($S_z$) states. Attitude, bias, accelerometer and magnetometer
models, the residual accelerometer-bias OU and its freeze, Joseph updates
and resets are OU-III's. \texttt{SeaStateFusionFilter\_CARMA3} keeps
OU-III's per-sample schedule (front end, proxy startup, staged commit,
watchdog, magnetic-service bias gate, accelerometer-noise inflation) and maps
the measured operating point onto every axis:
\begin{gather}
f_{\mathrm{tune}}=\frac{c_\tau}{2\tau},\qquad
\kappa=\frac{f_{\mathrm{tune}}}{f_{z,\mathrm{base}}},\nonumber\\
\sigma_z=\max\!\bigl(0.05,\ \sigma_{\mathrm{floor}},\ \sigma_{aw}/c_\sigma\bigr),
\qquad \sigma_h=S_{\mathrm{factor}}\,\sigma_z,
\label{eq:mapping}
\end{gather}
so the model's displacement zero-crossing frequency is placed on the
measured one and the amplitude is the measured band acceleration RMS with
OU-III's floor. No horizontal statistic is measured. The mappings were fixed
before any Monte Carlo run. $\mat{F}$ and $\mat{Q}$ are exact at a reference
step $h_0$, recomputed when the committed operating point changes or the step
moves more than $h_0/16$ from it; nearer steps take their Taylor series in
$h-h_0$ (\cref{tab:cost}).

\begin{table}[t]
\centering
\caption{State dimension and cost. Multiply-adds of one prediction and one
accelerometer update, and the multiply-adds a jittered step adds (Taylor
evaluation of $\mat{F}$, $\mat{Q}$ near the reference step).}
\label{tab:cost}
\footnotesize
\setlength{\tabcolsep}{3pt}
\begin{tabular}{@{}lcrr@{}}
\toprule
Configuration & $N_X$ & MACs & \makecell{jittered\\step}\\
\midrule
OU-III (shipping) & 21 & 12\,483 & 0\\
C41 $x/y/z$ & 21 & 4170 & 216\\
C52 $x/y/z$ & 24 & 5697 & 440\\
C41 $x/y$ + C52 $z$ & 22 & 4655 & 328\\
C52 $x/y$ + C41 $z$ & 23 & 5164 & 328\\
C52 $x/y/z$ + $S_x,S_y$ & 26 & 6971 & 610\\
\bottomrule
\end{tabular}
\end{table}

CARMA(4,1) everywhere has OU-III's 21 states and a third of its
multiply-adds per sample (\cref{tab:cost}): OU-III propagates its 12-state
linear block as one dense matrix on both sides of $\mat{P}$, while the CARMA
blocks are applied per axis to the upper block triangle, the gain is formed
through $\mat{S}^{-1}$ and the Joseph update is written as
$\mat{K}\mat{G}^\top-\mat{P}\mat{C}^\top\mat{K}^\top$ with
$\mat{G}=\mat{K}\mat{S}-\mat{P}\mat{C}^\top$ (the same expression for any
gain). A jittered sample clock costs only the Taylor evaluation while the
step stays within $h_0/16$ of the reference, at the float accuracy of the
exact discretization (host update at a fixed / $\pm1\,\%$ / $\pm25\,\%$
jittered interval: 2.53 / 2.60 / \SI{3.18}{\micro s} for C41 against 7.07 / 7.04 / \SI{7.12}{\micro s} for OU-III, whose
closed-form transition has no such cost). Host times rank the variants; they
are not ESP32-S3 times.

\section{Steady-State Analysis}
\label{sec:steady}
Exact steady-state Kalman transfers (algebraic Riccati equation at
\SI{200}{Hz}, white sensor noise, oracle operating point) applied to the true
displacement spectra separate signal distortion from noise propagation
(\cref{tab:steady}). OU-III's broad prior distorts moderately everywhere but
passes accelerometer noise at low frequency; a stationary CARMA prior
propagates two to four times less noise and wins whenever its shape is close
enough to the vessel spectrum. On the mixed swell and wind sea the fitted
shapes lose, CARMA(5,2) catastrophically: almost all of its error is
distortion.

\begin{table}[t]
\centering
\caption{Steady-state $x$ error, mm RMS (white noise, oracle operating
point); last row heave for OU-III / C41J / C52J.}
\label{tab:steady}
\footnotesize
\setlength{\tabcolsep}{3pt}
\begin{tabular}{@{}lrrrr@{}}
\toprule
Prior & JONSWAP & PM & mixed & bimodal\\
\midrule
OU-III & 10.0 & 13.7 & 70.2 & 33.2\\
C30J & 18.4 & 28.0 & 113.7 & 58.7\\
C41J & 5.6 & 8.8 & 87.6 & 23.3\\
C52J & 6.8 & 12.7 & 283.7 & 46.1\\
C41, sway-fitted & 4.8 & 7.5 & 80.4 & 23.2\\
\midrule
heave & \makecell{13.4/10.2\\/13.7} & \makecell{17.2/14.8\\/21.3} &
\makecell{60.7/93.7\\/302.1} & \makecell{31.5/26.9\\/51.3}\\
\bottomrule
\end{tabular}
\end{table}

\section{Paired Benchmark}
\label{sec:benchmark}
\texttt{tests/kalman\_carma/carma3\_bench} runs one scenario and seed per
process: a random-phase directional sea through the RAOs of
\cref{sec:physics} plus the scenario's low-frequency translation, sensed by
one IMU and magnetometer stream (\SI{200}{Hz}; accelerometer noise
\SI{1.5}{mg}, turn-on bias up to \SI{8}{mg} with a random walk; gyro bias
up to \SI{0.05}{\degree/s}; magnetometer at \SI{25}{Hz} with hard- and
soft-iron terms) that is fed identically to every filter. Errors are taken
against the wave coordinate over $t\ge\SI{300}{s}$ of a \SI{1200}{s} run.
The truth contains no CARMA model, and no ground truth enters any estimator
or tuner. The complete-INS level runs the filters as built; a level with
prescribed attitude isolates the motion prior. In total 14\,930 filter runs
were made; every table, confidence interval and Wilcoxon test is in
\texttt{reports/results/carma3\_benchmark/summary.md}.

\begin{table*}[t]
\centering
\caption{Complete INS, simulation: RMS error in \% of incident $H_s$,
horizontal $\sqrt{x^2+y^2}$ (H) and heave (Z); 100 paired seeds per sea;
still IMU in mm (30 seeds); NEES of position and velocity (6 if
consistent).}
\label{tab:principal}
\footnotesize
\setlength{\tabcolsep}{4pt}
\begin{tabular}{@{}lccccccccccc@{}}
\toprule
& & & \multicolumn{2}{c}{JONSWAP 0.29\,Hz} & \multicolumn{2}{c}{PM 0.25\,Hz} &
\multicolumn{2}{c}{Swell + wind sea} & \multicolumn{2}{c}{Still IMU (mm)} & NEES\\
Prior ($x/y$; $z$) & States & MACs & H & Z & H & Z & H & Z & H & Z & JONSWAP / mixed\\
\midrule
OU-III; OU-III & 21 & 12\,483 & 5.18 & 3.12 & 5.58 & 3.26 & 18.26 & 10.27 & 53 & 27 & 7.7 / 31\\
C41; C41 & 21 & 4170 & 2.66 & 1.83 & 2.91 & 2.10 & 15.14 & 11.65 & 84 & 53 & 5.4 / 64\\
C52; C52 & 24 & 5697 & 2.05 & 2.04 & 2.68 & 2.68 & 31.67 & 27.23 & 87 & 46 & 4.4 / 324\\
C52; C41 & 23 & 5164 & 2.05 & 1.83 & 2.67 & 2.10 & 31.08 & 11.65 & 87 & 53 & 4.1 / 190\\
C41; C52 & 22 & 4655 & 2.66 & 2.04 & 2.91 & 2.69 & 15.31 & 27.27 & 84 & 46 & 5.7 / 186\\
C52+$S$; C52 & 26 & 6971 & 2.02 & 2.04 & 2.77 & 2.68 & 31.76 & 27.23 & 80 & 46 & 4.5 / 336\\
\bottomrule
\end{tabular}
\end{table*}

\emph{Principal seas} (\cref{tab:principal}). CARMA(4,1) on all axes halves
the horizontal error on the single-peaked seas (paired change
$+48.7\,\%$, 95\,\% bootstrap interval $[+48.4,+48.9]$, on JONSWAP) and
lowers heave by 36--41\,\%, better than OU-III on every seed. On the mixed
swell and wind sea it is still 17\,\% better horizontally but 13\,\% worse in
heave. Heave depends only on the vertical model: variants that share it
differ by at most 0.04\,\% $H_s$ on any condition.

\begin{table}[t]
\centering
\caption{Further complete-INS conditions (30 seeds; 20 of \SI{3600}{s}):
RMS error in \% of $H_s$ for OU-III / C41 / C52 on all axes.}
\label{tab:conditions}
\footnotesize
\setlength{\tabcolsep}{2.5pt}
\begin{adjustbox}{max width=\columnwidth}
\begin{tabular}{@{}lcc@{}}
\toprule
Condition & horizontal & heave\\
\midrule
following, underway \SI{2.5}{m/s} & 10.54 / 4.03 / 4.04 & 8.43 / 4.72 / 5.25\\
head, underway \SI{2.5}{m/s} & 3.65 / 2.06 / 1.91 & 2.46 / 1.79 / 2.41\\
anchored slow drift & 6.67 / 3.64 / 2.64 & 3.16 / 1.82 / 2.04\\
drifting & 10.04 / 5.39 / 3.39 & 3.14 / 1.83 / 2.05\\
manoeuvring & 26.04 / 14.11 / 8.01 & 3.20 / 1.83 / 2.04\\
bimodal & 11.10 / 6.22 / 9.18 & 5.90 / 4.18 / 7.89\\
broadband & 16.85 / 7.34 / 7.60 & 8.32 / 5.39 / 5.79\\
JONSWAP \SI{0.10}{Hz}, \SI{4}{m} & 10.26 / 4.77 / 4.34 & 4.69 / 2.60 / 2.83\\
\SI{5}{cm} sea & 32.66 / 8.92 / 8.13 & 17.14 / 9.58 / 9.73\\
\SI{300}{s} magnetometer gap & 5.27 / 2.71 / 2.10 & 3.14 / 1.83 / 2.05\\
$\pm25\,\%$ $dt$ jitter & 5.26 / 2.71 / 2.10 & 3.15 / 1.84 / 2.07\\
mixed sea, \SI{3600}{s} & 18.30 / 15.09 / 31.72 & 10.25 / 11.61 / 27.06\\
\bottomrule
\end{tabular}
\end{adjustbox}
\end{table}

\emph{All conditions} (\cref{tab:conditions}). Horizontal CARMA(4,1) is
better than OU-III on all 31 complete-INS conditions: 43--73\,\% on every
single-peaked, bimodal, broadband, directional, underway, anchored, drifting,
manoeuvring, low-amplitude, sensor-fault and long case, 17--18\,\% on the
mixed sea. CARMA(5,2) is better than CARMA(4,1) on most single-peaked seas
and much better under low-frequency translation, worse on the bimodal sea,
and 73\,\% worse than OU-III on the mixed sea. With the attitude prescribed,
the ranking is unchanged, so the CARMA advantage is not an artifact of the
attitude coupling.

\emph{The mixed-sea collapse of CARMA(5,2).} Its numerator zeros sit at
$0.151\,\kappa$\,Hz. On the mixed sea the tuner's operating frequency has a
median of \SI{0.218}{Hz} ($\kappa\approx0.59$), which places the notch at
\SI{0.089}{Hz}, on the \SI{0.09}{Hz} swell carrying about 60\,\% of the
elevation variance. The swell then cannot be represented: 323 of 352\,mm of
$x$ error (seed 1) lies in \SIrange{0.05}{0.1}{Hz}, and the unexplained swell
acceleration leaks into the tilt (pitch \SI{1.6}{\degree} RMS). The mechanism
is structural, not a tuning issue; mixed orders inherit it.

\emph{Horizontal $S$.} CARMA(5,2) with the derived $S$ law is within a few
per cent of CARMA(5,2) everywhere and does not repair the mixed sea. Carrying
OU-III's $r_S$ onto CARMA(4,1), which Proposition~1 shows to be invalid,
lowers low-frequency errors (still IMU \SI{42}{mm}) but is worse where the sea
has broad or low-frequency energy: it acts as extra shrinkage that counts
prior information twice.

\emph{Still IMU.} Every CARMA prior shows more phantom displacement than
OU-III, horizontally (76--87 against \SI{53}{mm}) and in heave (46--53
against \SI{27}{mm}), nearly all below \SI{0.05}{Hz} and unchanged by
calibrating the turn-on bias. With no waves the tuner pins the operating
frequency at its \SI{0.05}{Hz} bound and the amplitude at its
\SI{0.05}{m/s^2} floor, so the resonant prior expects decimetres of
displacement at \SI{20}{s} periods and turns sensor noise and slow tilt and
bias leakage into motion there. In the first \SI{120}{s} after Live the
still-IMU heave transient is 2--3 times OU-III's; in waves it is smaller than
OU-III's, because CARMA starts from a stationary covariance that matches the
sea.

\emph{Consistency and attitude.} On single-peaked seas CARMA is closer to
consistent than OU-III, whose horizontal axes are overconfident; on the
mixed sea every filter is overconfident and CARMA much more
(\cref{tab:principal}): with a prior outside the true family the Kalman
covariance is not an error bound, and no model-error term is carried. Roll
and pitch errors are similar for every filter except CARMA(5,2) on the mixed
sea; the horizontal CARMA prior improves heading in the magnetometer gap
(\SI{2.09}{\degree} against \SI{2.90}{\degree}) and while manoeuvring.

\section{Pinned Vessel-RAO Dataset}
\label{sec:dataset}
The eight scored records of the dataset described above (JONSWAP and
PM--Stokes, $H_s$ \SIrange{0.27}{8.5}{m}) were replayed through the shared
\texttt{W3dSimulationRunner}: OU-III is the unchanged
\texttt{kalman\_ou\_iii-sim}, the CARMA variants \texttt{carma3-sim} with the
identical adapter configuration, ten seeds, final \SI{900}{s} of each record.
The hull differs from the synthetic benchmark's.

\begin{table}[t]
\centering
\caption{Pinned dataset: 10 seeds $\times$ 8 records, paired per (record,
seed); RMS in \% of $H_s$, change against OU-III with its 95\,\% bootstrap
interval (better on all 80 pairs).}
\label{tab:dataset}
\footnotesize
\setlength{\tabcolsep}{2.5pt}
\begin{adjustbox}{max width=\columnwidth}
\begin{tabular}{@{}lccc@{}}
\toprule
Filter & horizontal & heave & \makecell{roll/pitch/yaw\\(deg)}\\
\midrule
OU-III & 6.36 & 3.98 & 0.30/0.26/2.43\\
C41 $x/y/z$ & 3.04 (+52.2 [50.8, 53.8]) & 2.50 (+37.2 [35.3, 39.1]) & 0.30/0.24/2.29\\
C52 $x/y/z$ & 2.76 (+56.7 [54.5, 58.8]) & 2.83 (+28.9 [26.0, 31.7]) & 0.30/0.24/2.33\\
C52 $x/y$ + C41 $z$ & 2.76 (+56.7 [54.5, 58.8]) & 2.50 (+37.2 [35.3, 39.1]) & 0.29/0.24/2.33\\
\bottomrule
\end{tabular}
\end{adjustbox}
\end{table}

On this independent hull and these single-peaked seas the synthetic results
carry over (\cref{tab:dataset}). The dataset has no mixed sea, no still
water and no low-frequency translation, so it cannot test the obstacles of
\cref{sec:obstacles}.

\section{Deterministic Replay}
\label{sec:replay}
\Cref{tab:carma_sim_results} is generated during the build by
\texttt{tools/carma\_sim\_table.py} from \texttt{carma3-sim} (CARMA(4,1) on
all axes) with the default seeds, on the same eight records in the same
order as the OU-III and TFG tables. It is a single realization per record,
reported for continuity, not a controlled comparison; the controlled one is
\cref{tab:dataset}. Yaw in particular depends on the seed: the per-seed mean
yaw RMS of CARMA(4,1) over the eight records ranges from
\SIrange{0.95}{3.61}{\degree} across the ten seeds of \cref{tab:dataset},
and the default seeds give \SIrange{0.46}{0.75}{\degree} per record, below
all ten. The simulator fails the build if any
default-seed record exceeds its regression bar (worst value plus 0.5\,\%).

\IfFileExists{carma-sim-results-generated.tex-part}{%
  \input{carma-sim-results-generated.tex-part}%
}{%
  \begin{table}[t]\centering\small
    \fbox{\parbox{0.9\linewidth}{\centering Results table not generated.\\
      Run \texttt{tools/carma\_sim\_table.py}.}}
  \end{table}%
}

\begin{figure}[t]
  \centering
  \IncludePGF{w3d_carma3_jonswap_medium_xykin.pgf}
  \caption{Horizontal displacement, velocity and acceleration, JONSWAP
    $H_s=\SI{1.5}{m}$, last \SI{60}{s}; RMS errors over the scored
    \SI{900}{s} window.}
  \label{fig:xykin}
\end{figure}

\begin{figure}[t]
  \centering
  \IncludePGF{w3d_carma3_jonswap_medium_zkin.pgf}
  \caption{Vertical kinematics on the same record, with the
    $\pm$ displacement and velocity envelopes of the adaptation.}
  \label{fig:zkin}
\end{figure}

\begin{figure}[t]
  \centering
  \IncludePGF{w3d_carma3_jonswap_medium_tuner.pgf}
  \caption{Operating point over the whole replay: tracker frequency and
    $1/(2\tau_{\mathrm{applied}})$, the frequency on which the CARMA shapes
    are placed; band acceleration and the applied $\sigma_a$; $\tau$; and the
    vertical scale $\kappa_z$.}
  \label{fig:tuner}
\end{figure}

\section{Obstacles and Recommendation}
\label{sec:obstacles}
On this evidence CARMA(4,1) on all three axes is the research candidate. It
is preferred to OU-III's horizontal chain with CARMA in heave, because
horizontal CARMA(4,1) wins on all 31 conditions at no extra cost and heave
does not depend on the horizontal prior; to CARMA(5,2), whose notch the
frequency mapping can place on a swell; to mixed orders, which inherit that
failure; and to coupled priors, whose oracle gain is about 5\,\% and for which
no estimator of the coupling exists. It is \emph{not} shippable. Remaining
obstacles:
\begin{enumerate}[leftmargin=*,itemsep=1pt]
\item still water and very quiet seas: phantom displacement and a still-IMU
  heave startup transient 2--3 times OU-III's; a valid low-frequency or
  ``no waves'' mechanism in the prior is needed;
\item the mixed swell and wind sea: heave 13\,\% worse than OU-III; one fixed
  shape scaled by one frequency cannot represent two peaks;
\item covariance under mismatch: no model-error term, so the covariance is
  not an error bound on multi-peaked seas;
\item hull and sea dependence: one synthetic and one replay hull, a
  JONSWAP-fitted horizontal shape, no lever arm, wind heel or nonlinear
  waves;
\item no hardware evidence (AtomS3R timing, float behaviour, sea data);
\item no stability analysis.
\end{enumerate}

\section{Relation to the OU-III Stability Study}
The OU-III stability study and its premises (marine motion, IMU bias and
magnetic service, applied to one persistent physical execution) concern the
shipping OU-III filter. None of them transfers to this filter. A
nondegenerate stationary Gaussian CARMA realization is not almost surely
bounded over an infinite horizon, and its jerk is not subject to a
deterministic bound; the operating-point commits make the filter a switched
linear system. A replacement of OU-III would require its own process-model
mismatch, observability, covariance and full estimator stability analyses.

\section*{Reproducibility Note}
The filter, its wrapper and the CARMA model library are in
\texttt{src/kalman\_carma}; the tests, the benchmark and the dataset
simulator are in \texttt{tests/kalman\_carma}. The benchmark driver
\texttt{tools/carma3\_benchmark.py} reproduces every run from committed
seeds, and its raw results are in
\texttt{reports/results/carma3\_benchmark}. This paper and its chart gallery
are built by the repository's GitHub Actions workflow from a fresh replay of
the dataset.

\begin{thebibliography}{9}
\bibitem{brockwell01}
P. J. Brockwell, ``Continuous-time ARMA processes,'' in \emph{Handbook of
Statistics}, vol. 19, Elsevier, 2001, pp. 249--276.

\bibitem{brockwell04}
P. J. Brockwell, ``Representations of continuous-time ARMA processes,''
\emph{J. Appl. Probab.}, vol. 41, no. A, pp. 375--382, 2004.

\bibitem{sarkka23}
S. S\"arkk\"a and L. Svensson, \emph{Bayesian Filtering and Smoothing},
2nd ed., Cambridge University Press, 2023.
\end{thebibliography}

\end{document}
